\documentclass[10pt,a4paper]{amsart}
\usepackage[margin=19mm]{geometry}
\usepackage{amsmath,amssymb,mathtools}
\usepackage[scr=rsfs,cal=euler]{mathalfa}
\usepackage{xfrac}
\usepackage[unicode,hidelinks,bookmarks=false]{hyperref}
\newcommand{\ie}{i.e.\ }
\newcommand{\cf}{cf.\ }
\newcommand{\resp}{respectively\ }
\newcommand{\Subsection}{\subsection}
\providecommand{\nobreakdash}{\nobreak\hbox{-}}
\setlength{\emergencystretch}{2em}
\multlinegap0pt
\usepackage{bm}
\usepackage{bbm}
\usepackage{dsfont}
\usepackage{xspace}
\usepackage{cancel}
\usepackage{mathtools}
\usepackage{amsthm}
\makeatletter
\@ifundefined{thm}{\newtheorem{thm}{Thm}}{}
\@ifundefined{cor}{\newtheorem{cor}[thm]{Corollary}}{}
\@ifundefined{lemma}{\newtheorem{lemma}[thm]{Lemma}}{}
\@ifundefined{prop}{\newtheorem{prop}[thm]{Proposition}}{}
\makeatother
\newcommand{\eps}{\varepsilon}
\title[The Quadratic and Cubic Characters of 2]{The Quadratic and Cubic Characters of 2}
\author{Matias C. Relyea}
\date{Source: arXiv:2410.21646v3 (CC BY 4.0)}
\begin{document}
\maketitle

\noindent\emph{Source and licence.} The Quadratic and Cubic Characters of 2. Version arXiv:2410.21646v3 (CC BY 4.0), distributed under the Creative Commons Attribution 4.0 International licence, as recorded in the arXiv licence metadata for this exact version and shown on the abstract page https://arxiv.org/abs/2410.21646v3. Full rights notice: licenses/arxiv-2410.21646-CC-BY-4.0.txt.\par\smallskip

\noindent\textit{Editorial scope.} Counted unit: the Lemma whose source label is \texttt{Lemma: value of the general Gauss sum} in the article (source0.tex lines 814--816, source label \texttt{Lemma: value of the general Gauss sum}), delivered together with the article's own complete original proof of it (source0.tex lines 817--850, one \texttt{proof} environment of 34 lines). No mathematics is written, repaired or re-derived: statement and proof are the article's own text line for line. Required context retained uncounted: Proposition: proof about sum of elements evaluated by a character prop (lines 696--698); Lemma: two properties of the Gauss sum (lines 725--731); Lemma: result about zeta from special property prop about Gauss sums about zeta from special property prop about Gauss sums (lines 751--759); corollary to Lemma 3.7 (lines 763--767). Every definition, display and label of those lines is retained unchanged. Omitted from this package and not redistributed: 1--695, 699--724, 732--750, 760--762, 768--813, 851--1125. Nothing inside a retained excerpt is omitted or altered: every retained body is delivered exactly as the article's lines are written, so no omission receipt of any kind accompanies this package. Citations: the retained excerpts carry no citation command at all, so no bibliography entry of the article is transcribed and no reference is redistributed. Reference adaptation: none was needed and none was made; every reference inside the delivered unit resolves inside it, so no unresolved reference or citation is produced. Printed slips kept literal: no printed word, symbol, hypothesis or number of the counted statement or of its proof was altered. Layout and class: the article's own class is \texttt{amsart} with packages including etex, tabularx, fancyhdr, float, geometry, babel, inputenc, subcaption, amsmath, amssymb, amsfonts, amsthm, mathdots, mathrsfs; the wrapper is the lane's pinned \texttt{amsart} editorial wrapper at 10pt on A4, which restates the theorem-like environments the retained text needs and adds the article's own macro definitions that the retained text needs (\texttt{\textbackslash eps}), with extra wrapper packages (bm, bbm, dsfont, xspace, cancel, mathtools). The delivered numbering is the unit's own. Assets: no retained span contains a figure environment, a graphics-inclusion command, a PDF-page inclusion, a table or a TikZ picture; no image file is copied into this package, the package declares no asset dependency and no image file is redistributed with it. Licence: version arXiv:2410.21646v3 of this article is distributed under the Creative Commons Attribution 4.0 International licence, as recorded in the arXiv licence metadata for this exact version and shown on the abstract page https://arxiv.org/abs/2410.21646v3. A full rights notice is delivered at licenses/arxiv-2410.21646-CC-BY-4.0.txt. Characters: every retained excerpt is pure ASCII, so the delivered engine, XeLaTeX with its default Latin Modern text font, reports no missing character.
\begin{prop} \label{Proposition: proof about sum of elements evaluated by a character prop}
    Let $\chi$ be a multiplicative character. If $\chi\neq\eps$, the trivial multiplicative character, then $\sum_{t\in \mathbb{F}_{p}}\chi(t)=0$. Otherwise, the sum is $p$. 
\end{prop}

\begin{lemma} \label{Lemma: two properties of the Gauss sum}
The following are true. 
    \begin{enumerate}
        \item If $a\neq 0$ and $\chi\neq \eps$, then $g_{a}(\chi)=\overline{\chi(a)}g_{1}(\chi)$. 
        \item If $a\neq 0$ and $\chi = \eps$, then $g_{a}(\eps)=0$.
    \end{enumerate}
\end{lemma}

\begin{lemma} \label{Lemma: result about zeta from special property prop about Gauss sums about zeta from special property prop about Gauss sums}
    \begin{equation*}
        \sum_{t=0}^{p-1}\zeta_{p}^{at} = \left\{ \begin{array}{ll}
            p, & a\equiv 0\pmod{p}, \\
            0, & a\not\equiv 0\pmod{p}.
        \end{array}
        \right.
    \end{equation*}
\end{lemma}

\begin{cor}[Corollary to Lemma \ref{Lemma: result about zeta from special property prop about Gauss sums about zeta from special property prop about Gauss sums}]\label{corollary to Lemma 3.7}
    \begin{equation*}
     p^{-1}\sum_{t=0}^{p-1}\zeta_{p}^{t(x-y)}  = \delta(x,y).
\end{equation*}
\end{cor}

\begin{lemma}\label{Lemma: value of the general Gauss sum}
    If $\chi\neq \varepsilon$ is a nontrivial character, then $|g(\chi)|^{2}=p$. 
\end{lemma}

\begin{proof}
    The main idea for the proof is to evaluate the sum 
    \begin{equation*}
        \sum_{a\in\mathbb{F}_{p}}g_{a}(\chi)\overline{g_{a}(\chi)}
    \end{equation*}
    in two different ways and set the evaluations equal to one another. We will first evaluate the argument contained within the sum. Assume that $a\neq 0$. By (1) of Lemma \ref{Lemma: two properties of the Gauss sum}, we can write 
    \begin{equation*}
        \overline{g_{a}(\chi)} = \overline{\chi(a^{-1})g(\chi)} = \chi(a)\overline{g(\chi)}.
    \end{equation*}
    Taking the conjugate, we also have $g_{a}(\chi) = \chi(a^{-1})g(\chi)$. Multiplying and rearranging, we have
    \begin{equation*}
    \chi(a)\overline{g(\chi)}\chi(a^{-1})g(\chi) = \overline{g(\chi)}g(\chi) = |g(\chi)|^{2}.
    \end{equation*}
    Since $\sum_{a\in\mathbb{F}_{p}}$ sums over all elements of $\mathbb{F}_{p}$ except $a=0$, we consider this quantity $p-1$ times. So,
    \begin{equation*}
        \sum_{a\in\mathbb{F}_{p}}g_{a}(\chi)\overline{g_{a}(\chi)} = (p-1)|g(\chi)|^{2}.
    \end{equation*}
    Similarly, considering two parameters $x$ and $y$ and rewriting the argument as a double sum, we have
    \begin{equation*}
        g_{a}(\chi)\overline{g_{a}(\chi)} = \sum_{x\in\mathbb{F}_{p}}\sum_{y\in\mathbb{F}_{p}}\chi(x)\zeta_{p}^{ax}\overline{\chi(y)}\zeta_{p}^{ay} = \sum_{x\in\mathbb{F}_{p}}\sum_{y\in\mathbb{F}_{p}}\chi(x)\overline{\chi(y)}\zeta_{p}^{ax-ay}.
    \end{equation*}
    Summing over all elements of $\mathbb{F}_{p}$ and applying Corollary \ref{corollary to Lemma 3.7}, we have
    \begin{align*}  \sum_{a\in\mathbb{F}_{p}}\sum_{x\in\mathbb{F}_{p}}\sum_{y\in\mathbb{F}_{p}}\chi(x)\overline{\chi(y)}\zeta_{p}^{ax-ay} & = pp^{-1}\sum_{a\in\mathbb{F}_{p}}\sum_{x\in\mathbb{F}_{p}}\sum_{y\in\mathbb{F}_{p}}\chi(x)\overline{\chi(y)}\zeta_{p}^{ax-ay}p \\ & = p\sum_{x\in\mathbb{F}_{p}}\sum_{y\in\mathbb{F}_{p}}\chi(x)\overline{\chi(y)}\zeta_{p}^{ax-ay}\delta(x,y).
    \end{align*}
    If $x\not\equiv y\pmod{p}$ then the double sum will equate to $0$ as $\delta(x,y)$ will be $0$, and by Proposition \ref{Proposition: proof about sum of elements evaluated by a character prop} the sum of all multiplicative characters over $\mathbb{F}_{p}$ is 0. Therefore we consider when $x\equiv y\pmod{p}$. If this is true, then every term except when $x\equiv y\pmod{p}$ will be counted, leaving a total of $p-1$ terms. Since $x\equiv y\pmod{p}$, necessarily $\zeta_{p}^{ax-ay}=1$, so
    \begin{equation*}
        p\sum_{x\in\mathbb{F}_{p}}\sum_{y\in\mathbb{F}_{p}}\chi(x)\overline{\chi(y)}\zeta_{p}^{ax-ay}\delta(x,y) = p(p-1).
    \end{equation*}
    Equating our two evaluations, we have 
    \begin{align*}
        (p-1)|g(\chi)|^{2} & = p(p-1) \\
        |g(\chi)|^{2} & = p.
    \end{align*}
\end{proof}

\end{document}
